\begin{lemma}[Graded ranks and fixed simples]
\label{lem:m3-graded-rank-untwisting}
In a faithful finite-group graded complex fusion category with neutral
component \(\mathcal D\) and \(\operatorname{BrPic}(\mathcal D)=1\),
every homogeneous component has rank \(\operatorname{rank}(\mathcal D)\).
If a pointed subcategory \(\operatorname{Vec}_G^\omega\) fixes a simple
object, then \([\omega]=0\).  Both statements hold without unitarity.
\end{lemma}
\begin{proof}
Each homogeneous component is an invertible \(\mathcal D\)-bimodule
\cite[Theorem~6.1(i)]{ENOExt}, hence belongs to the regular class;
equivalence preserves rank.  For the second assertion, the fixed
simple spans a rank-one module category.  Choosing isomorphisms for
the group action gives scalar composition factors, and the module
pentagon writes \(\omega\) as their coboundary.
\end{proof}

\begin{proposition}[The Haagerup graded-component obstruction]
\label{prop:m3-r455-unitary}
The ring \(R8\_3\_455\) has no unitary categorification.  The argument
does not exclude general complex categorification.
\end{proposition}
\begin{proof}
The source basis splits into neutral indices \(\{0,1,2,3,6,7\}\) and
odd indices \(\{4,5\}\), giving a faithful \(C_2\)-grading.  The neutral
ring has Haagerup rules
\[
 g^3=1,\qquad \rho g=g^{-1}\rho,\qquad
 \rho^2=1+\rho+g\rho+g^2\rho.
\]
The simple \(Y\) at source index \(4\) satisfies \(gY=Yg=Y\) and
\(Y^2=1+g+g^2\).  The certificate independently reconstructs the full
tensor and checks all \(512\) coefficients and the grading.

Let \(\mathcal D\) be the neutral category in a hypothetical unitary
categorification.  The additive category generated by \(Y\) is a
rank-one module over its pointed subcategory
\(\operatorname{Vec}_{C_3}^{\omega}\).  Choosing isomorphisms
\(g\otimes Y\cong Y\), the module pentagon expresses \(\omega\) as
the coboundary of their composition scalars.  Hence
\(c^{0,3}(\mathcal D)=0\).  Also
\(H^2(C_3,\mathbb C^\times)=0\), so Izumi's second obstruction
\(c^{1,2}(\mathcal D)\in H^1(C_2,H^2(C_3,\mathbb C^\times))\)
vanishes.  These steps establish the hypotheses of
\cite[Definition~2.7]{Izumi3n}; omitting the first step would leave the
twisted alternatives mentioned immediately after that definition.

By \cite[Theorems~5.9 and 9.1]{Izumi3n}, the unitary category
\(\mathcal D\) is one of the two known categories \(H_2,H_3\).
Their Brauer--Picard groups are trivial
\cite[Theorem~1.5]{GSHaagerup}.  Every homogeneous component of a
faithful extension is an invertible bimodule over its neutral category
\cite[Theorem~6.1(i)]{ENOExt}.  The odd component must consequently be
equivalent to the regular \(\mathcal D\)-bimodule, whose rank is six.
Its rank is two, a contradiction.  The classification invoked here is
a C*-classification and supplies no such conclusion for hypothetical
nonunitary neutral categories.
\end{proof}
